\clearpage\begingroup
\hypertarget{native-338b}{}\HistoricalHead{333}{R}{BEOBACHTUNGEN AM REICHENBACHSCHEN MERIDIANKREISE.}
\begin{center}\textsc{Gauss} an \textsc{Olbers}, 1. Februar 1824.\end{center}
\noindent\hspace*{\parindent}. . . . . . Sehr befremdet hat mich die Stelle Ihres vorletzten Briefes, dass \textsc{Bessel} geäussert haben soll, er halte meine Polhöhe für \(1\overset{\prime\prime}{,}5\) zu gross. Sie drücken Sich so aus, als halten Sie für zweifelhaft, ob er dieses wirklich geäussert habe, und nehmen es nur unter der Voraussetzung an, dass ich seine Refraction aus Gründen nicht gebraucht hätte. Allein diese Voraussetzung hat \textls[120]{nicht} Statt, und \textsc{Bessel} \textls[120]{weiss}, dass ich seine Refraction
\clearpage\SourceLeaf{339}\HistoricalHead{334}{L}{PRAKTISCHE UND SPHÄRISCHE ASTRONOMIE. BRIEFWECHSEL.}
\noindent gebraucht habe, obgleich seine Gründe mich keinesweges überzeugt haben, dass es nothwendig sei, die \textls[120]{äussere} Temperatur zum Grunde zu legen. Letzteres habe ich jedoch bisher der Gleichförmigkeit wegen wirklich gethan.

Was übrigens meine Polhöhe betrifft, so ist sie allerdings insofern hypothetisch, als ich die Inflexion des Fernrohrs noch nicht hinreichend untersucht habe. Insofern diese Inflexion vernachlässigt wird, kann ich von meiner Polhöhe nach den Beobachtungen von 1820 \textls[120]{Nichts} ablassen. Die eigne Untersuchung der Inflexion halte ich keinesweges für unnöthig, sondern sie ist nur verschoben, bis die \textls[120]{Zeit} dazu sein wird. Dass ich sie aber bisher \(=0\) gesetzt habe, geschah
\par\noindent\hspace*{\parindent}1) weil die wenigen Beobachtungen aus dem Wasserhorizont gar keine merkliche Flexion anzeigen,
\par\noindent\hspace*{\parindent}2) weil \textsc{Bessels} Declinationen der südlichen Sterne (Fundamentalsterne) sehr gut mit den meinigen übereinstimmen.

Es ist an \textls[120]{sich} nicht unmöglich, dass zahlreiche Beobachtungen eine nicht ganz unmerkliche Flexion geben u[nd] vielleicht meine Polhöhe um \(1''\) kleiner geben können. Allein \textls[120]{dann} würden \textls[120]{alle meine} Bestimmungen der Fundamentalsterne südlicher werden als die \textsc{Bessel}schen u[nd] ich habe durchaus keinen Grund, meine Bestimmungen für weniger genau zu halten als die \textsc{Bessel}schen.

Auf keinen Fall kann ich darin auch nur den kleinsten Grund zur Verminderung meiner Polhöhe finden, dass ihre Uebertragung auf München vermittelst der geodätischen Messungen \(1\overset{\prime\prime}{,}7\) mehr gibt als die dortigen Messungen. . . . . . .

Wenn man übrigens überlegt, was für eine ungeheure Arbeit erfordert wird, um eine Ungewissheit von \(1''\) auf eine Ungewissheit von \(\tfrac{1}{2}''\) zu reduciren, und zwar eine an sich durchaus geistlose Arbeit, so wird man zweifelhaft, ob der Zweck so viel werth ist wie die Mittel, und ob man es verantworten könne, einen grossen Theil der Kräfte eines Menschenlebens darauf zu wenden. . . .

\begin{center}\textsc{Gauss} an \textsc{Olbers}, 2. Februar 1824.\end{center}
\noindent\hspace*{\parindent}Nothwendig . . . . . muss ich heute noch einmal auf die Polhöhe meiner Sternwarte zurückkommen.
\clearpage\SourceLeaf{340}\HistoricalHead{335}{R}{BEOBACHTUNGEN AM REICHENBACHSCHEN MERIDIANKREISE.}
Ich habe Ihnen gestern geschrieben, dass ich, 1) \textls[120]{wenn mein Instrument keine Flexion hat}, von meiner Polhöhe \(51^0\,31'\,48\overset{\prime\prime}{,}7\) nichts ablassen kann, 2) dass ich das Dasein einer merklichen Biegung an sich gar nicht für unmöglich halte, da die sehr wenigen bisherigen Beobachtungen aus dem Wasserhorizont (die genau genommen eine \textls[120]{kleine} Biegung in dem Sinne, dass das Instrument die Z[enith]d[istanzen] zu klein angibt, wirklich andeuteten) lange nicht zureichen, dies zu entscheiden, 3) dass aber, wenn eine Biegung (in dem Sinne, wie ich gesagt habe, eine \textls[120]{negative} Biegung) Statt fände, sie die Uebereinstimmung meiner Declinationen der Fundamentalsterne mit den \textsc{Bessel}schen zerstören u[nd] meine südlicher geben würde, und dass ich aus \textls[120]{dieser} Ursache bisher immer geglaubt habe, die Biegung könne nicht merklich sein, da ich mir bewusst bin, meine eignen Beobachtungen mit Sorgfalt u[nd] schöner Uebereinstimmung unter sich gemacht zu haben.

Sie sehen, die Eine Hälfte meines Arguments beruht gerade auf dem Vertrauen selbst, was ich zu \textsc{Bessels} Bestimmungen habe.

Allein heute gestaltet sich die Sache ganz anders.

Meine Declinationen der Fundamentalsterne stimmten mit den \textsc{Bessel}schen sehr schön überein, ohne einen Ausschlag nach einer Seite zu zeigen, nach dem Verzeichniss der Declinationen, welches \textsc{Bessel} mir, vor 3 Jahren etwa, in einem Briefe[\begingroup\renewcommand{\thefootnote}{*)}\footnotemark\renewcommand{\thefootnote}{[*)}\footnotetext{In dem Briefe vom 19. April 1821, S. 333.]}\endgroup]\ schickte. \textls[120]{Dies} habe ich immer im Gedächtnisse gehabt und gar nicht daran gedacht, es mit dem später gedruckten zu vergleichen. Der 7. Band von \textsc{Bessels} Beob[achtungen] war sogar bisher noch nicht gebunden, da ich mich begnügt hatte, die Einleitung im allgemeinen zu durchlaufen, ohne die Zahlen mit den Briefen zu vergleichen.

Allein heute bekomme ich diesen Band vom Buchbinder zurück u[nd] sehe, dass diese Declinationen alle \textls[120]{ganz anders} sind als die im Briefe. Jetzt sind von 35 Fundamentalsternen (\(\alpha^2\) Capricorni habe ich gar nicht beobachtet) 34 bei mir nördlicher und nur einer, \(\alpha'\) Capricorni, den ich aber nur Einmal beobachtet habe, \(0\overset{\prime\prime}{,}16\) südlicher.

Mein Argument fällt also hienach ganz weg, und ich halte es nunmehro, \textls[120]{wieder im Vertrauen auf \textsc{Bessels} Bestimmungen}\begingroup\renewcommand{\thefootnote}{**)}\footnote{Denn meine durch Flexion nicht verbesserten Declinationen weichen von den \textsc{Bessel}schen in demselben Sinn ab, obwohl nicht so viel, wie die \textls[120]{aller andern} Astronomen \textsc{Piazzi}, \textsc{Brinkley}, \textsc{Oriani} u[nd] \textsc{Pond}.}\endgroup\ für \textls[120]{höchst}
\clearpage\SourceLeaf{341}\HistoricalHead{336}{L}{PRAKTISCHE UND SPHÄRISCHE ASTRONOMIE. BRIEFWECHSEL.}
\noindent\textls[120]{wahrscheinlich}, dass alle meine Zenithdistanzen zu klein sind und die Polhöhe zu gross.

Auch meine Circumpolarsterne, die alle ganz unabhängig von der Polhöhe durch obere und untere Culmination bestimmt sind, stimmen damit überein. Ich habe 1820 37 solcher Polarabstände bestimmt; unter diesen sind 17, die \textsc{Bessel} auch hat, und unter diesen sind 14, wobei meine Polardistanz kleiner ist als bei \textsc{Bessel}, wie eine negative Flexion meines Instruments es hervorbringen müsste; die drei, wobei meine Polardistanz \textls[120]{grösser} ist, sind
\begin{center}\begin{tabular}{@{}lr@{}}
Nordstern & \(+0\overset{\prime\prime}{,}16\)\\
\(\zeta\) Ursae min. & \(+0{,}95\)\\
\(\beta\) Ursae min. & \(+0{,}64\)
\end{tabular}\end{center}
\noindent alle dem Pol noch näher, wo also die Flexion jene Wirkung noch nicht so entschieden zeigen kann. Beim Nordstern ist ohnehin der Unterschied fast 0, und von \(\zeta\) Ursae min. habe ich nur 2 Beob[achtungen] (bei \(\beta\) das Gewicht von \(12{,}6\) Beob[achtungen]).

Ich werde einstweilen aus der Vergleichung mit \textsc{Bessels} Bestimmung die Flexion abzuleiten suchen[\begingroup\renewcommand{\thefootnote}{*)}\footnotemark\renewcommand{\thefootnote}{[*)}\footnotetext{Vgl. Notiz 5 und den Brief an \textsc{Bessel} vom 14. März 1824, S. 338.]}\endgroup], bis sich Zeit u[nd] Gelegenheit findet, sie direct selbst zu bestimmen.

Wenn ich die Declinationen der von mir im Aug[ust] 1820 beob[achteten] Zenithalsterne von \textsc{Bessel} entlehnen wollte, so würden diese Sterne folgendes Resultat meiner Polhöhe geben:
\begin{center}\begingroup\fontsize{9.5}{12}\selectfont\setlength{\tabcolsep}{5pt}
\begin{tabular}{@{}lr@{\hspace{15pt}}lr@{}}
\(\eta\) Cephei & \(51^0\,31'\,46\overset{\prime\prime}{,}82\) & P[\textsc{iazzi}] XX. 391 & \(51^0\,31'\,46\overset{\prime\prime}{,}73\)\\
P[\textsc{iazzi}] XX. 222 & \(47{,}56\) & 51 Drac[onis] & \(46{,}76\)\\
47 Drac[onis] & \(47{,}52\) & \(\kappa\) Cygni & \(46{,}82\)\\
Ceph[ei] 2 Hev. & \(47{,}65\) & P[\textsc{iazzi}] XXI. 32 & \(46{,}54\)\\
48 Drac[onis] & \(46{,}72\) & 20 Cygni & \(46{,}32\)\\
53 Dracon[is] & \(46{,}94\) & \(\iota\) Cygni & \(46{,}36\)\\
33 Cygni & \(46{,}60\) & \(c'\) Cygni pr[aec]. & \(47{,}01\)\\
49 Dracon[is] & \(46{,}77\) & \(\prime\prime\)\quad sequ. & \(46{,}70\)\\
46 Dracon[is] & \(47{,}46\) & \(\theta\) Cygni & \(46{,}17\)\\
 & & \(\omega'\) Cygni & \(47{,}53\)\\\hline
\multicolumn{3}{@{}l}{Mittel ohne Rücksicht auf Anz[ahl] d[er] B[eobachtungen]} & \(51^0\,31'\,46\overset{\prime\prime}{,}84\)
\end{tabular}\endgroup\end{center}
\clearpage\SourceLeaf{342}\HistoricalHead{337}{R}{BEOBACHTUNGEN AM REICHENBACHSCHEN MERIDIANKREISE.}\HistoricalFoot{XI1.}{43}
Wären \textsc{Bessels} Declinationen absolut genau, wofür natürlich diese schöne Uebereinstimmung an sich \textls[120]{gar nichts} beweisen kann, so würde der »\textls[120]{mittlere}« Fehler, welcher in diesem Endresultat zu befürchten ist, \(=0\overset{\prime\prime}{,}104\), oder der sogenannte wahrscheinliche Fehler \(=0\overset{\prime\prime}{,}07\) sein.

Die Münchener Polhöhe würde also ganz übereinstimmen. Ich wiederhole aber, dass ich \textls[120]{darin} durchaus \textls[120]{gar keinen} Grund für die absolute Richtigkeit finden kann. . . . . . .

\begin{center}\textsc{Gauss} an \textsc{Olbers}, 28. Februar 1824[\begingroup\renewcommand{\thefootnote}{*)}\footnotemark\renewcommand{\thefootnote}{[*)}\footnotetext{Vgl. Notiz 5, S. 348.]}\endgroup].\end{center}
\noindent\hspace*{\parindent}. . . . . . Wenn ich annehme, dass \textsc{Bessels} Bestimmung der Flexion des Fernrohrs seines Meridiankreises richtig ist, so wird die Flexion des Fernrohrs an dem meinigen aus der Vergleichung von 34 Fundamentalsternen, 1820 beobachtet,
\[
=1\overset{\prime\prime}{,}51\sin z,
\]
\noindent welche Grösse zu der beobachteten Zenithdistanz addirt werden muss. Für die Polhöhe finde ich dann aus sämmtlichen Circumpolarstern-Beobachtungen von 1820
\[
51^0\,31'\,47\overset{\prime\prime}{,}92\quad\text{oder}\quad51^0\,31'\,48\overset{\prime\prime}{,}03,
\]
\noindent je nachdem die Zahl der Beobachtungen in jeder Periode berücksichtigt oder nicht berücksichtigt wird. Ich nehme daher einstweilen \(51^0\,31'\,48\overset{\prime\prime}{,}0\) an, bis ich eine ganz selbstständige Bestimmung der Flexion machen kann. Man darf an diesem Resultate nichts erhebliches ändern, ohne die Uebereinstimmung der \textsc{Bessel}schen Beobachtungen mit den meinigen, 1820, zu vernichten. Da der \textsc{Repsold}sche Kreis jetzt doch nicht aufgestellt ist und auch wohl sobald nicht wieder aufgestellt werden kann, da der Zenithsector an diesem Platze steht, so lasse ich jetzt eine Einrichtung machen, um die zu jenem Kreise gehörenden schönen Mikroskope zur Prüfung der Theilung des \textsc{Reichenbach}schen Kreises zu verwenden. Inzwischen wird nun wohl vor künftigem Winter an den Gebrauch nicht zu denken sein. . . . . . .
\clearpage\endgroup
